\subsubsection{Data Integration}
The R module provides comprehensive routines for data integration, based on Jensen-Shannon Divergence. This is very useful for cross-study data.

\textbf{Key Features:}
\begin{itemize}
    \item \textbf{Calculate Jensen-Shannon Compatibility Score}: Score to evaluate how compatible two studies are.
    \item \textbf{Perform JSD Permutation Test}: Random sampling of the input data and recalculation of the JS score to test stability of the compatibility.
    \item \textbf{Family-Wise Score Contribution}: Calculates the impact of selected families on the calculated score.
\end{itemize}

\begin{enumerate}
    \item \textbf{Compute Gene Means}

        Compute per‑gene mean expression values across replicates, handling \lstinline!NaN! values appropriately.

        \textbf{Arguments:}
        \begin{itemize}
            \item \lstinline!expr!: numeric matrix of expression values \lstinline!(replicates * genes)!.
        \end{itemize}

        \textbf{Returns:}
        \begin{itemize}
            \item \lstinline!numeric!: vector of per‑gene mean expression values.
        \end{itemize}

        \begin{lstlisting}[language=R, caption={Per‑Gene Mean Expression}]
means <- tox_compute_gene_means(expr)
        \end{lstlisting}

    \item \textbf{Compute Signed Residuals}

        Compute signed residuals for each gene by subtracting the per‑gene mean expression from each replicate.

        \textbf{Arguments:}
        \begin{itemize}
            \item \lstinline!expr!: numeric matrix \lstinline!(replicates * genes)!.
            \item \lstinline!means!: numeric vector of per‑gene mean expression values.
        \end{itemize}

        \textbf{Returns:}
        \begin{itemize}
            \item \lstinline!numeric!: matrix of signed residuals \lstinline!(replicates * genes)!.
        \end{itemize}

        \begin{lstlisting}[language=R, caption={Signed Residual Computation}]
resid <- tox_compute_residuals(expr, means)
        \end{lstlisting}

    \item \textbf{Pool Mean Expression Values}

        Pool per‑gene mean expression values from two studies and compute reference points for neighborhood construction.

        \textbf{Arguments:}
        \begin{itemize}
            \item \lstinline!mean_S1!: numeric vector of mean expression values for study 1.
            \item \lstinline!mean_S2!: numeric vector of mean expression values for study 2.
            \item \lstinline!n_points!: integer number of reference points to compute.
        \end{itemize}

        \textbf{Returns:}
        \begin{itemize}
            \item \lstinline!list!: \{
                \lstinline!n_pool!: integer number of valid pooled means,\\
                \lstinline!x_star!: numeric vector of reference points
            \}
        \end{itemize}

        \begin{lstlisting}[language=R, caption={Pooling Mean Expression Values}]
out <- tox_pool_means(mean_S1, mean_S2, n_points)

n_pool <- out$n_pool
x_star <- out$x_star
        \end{lstlisting}

    \item \textbf{Pool Mean Expression Values (Expert)}

        Pool mean‑expression values using pre‑pooled and pre‑sorted arrays.  
        This expert interface bypasses internal pooling and sorting, allowing direct computation of reference points.

        \textbf{Arguments:}
        \begin{itemize}
            \item \lstinline!pooled_means!: numeric vector of pooled mean expression values.
            \item \lstinline!pooled_perm!: integer vector giving the sorting permutation for \lstinline!pooled_means!.
            \item \lstinline!n_points!: integer number of reference points to compute.
        \end{itemize}

        \textbf{Returns:}
        \begin{itemize}
            \item \lstinline!list!: \{
                \lstinline!n_pool!: integer number of valid pooled means,\\
                \lstinline!x_star!: numeric vector of reference points
            \}
        \end{itemize}

        \begin{lstlisting}[language=R, caption={Pooling Means (Expert Interface)}]
out <- tox_pool_means_expert(
  pooled_means,
  pooled_perm,
  n_points
)

n_pool <- out$n_pool
x_star <- out$x_star
        \end{lstlisting}

    \item \textbf{Construct Neighborhoods}

        Construct neighborhood‑based residual sets by selecting, for each reference point, the closest genes in mean‑expression space.  
        The neighborhood size is determined automatically unless a desired size is specified.

        \textbf{Arguments:}
        \begin{itemize}
            \item \lstinline!x_star!: numeric vector of reference mean‑expression points.
            \item \lstinline!n_pool!: integer number of pooled mean‑expression values.
            \item \lstinline!mean_S!: numeric vector of per‑gene mean expression values.
            \item \lstinline!resid_S!: numeric matrix of signed residuals \lstinline!(replicates * genes)!.
            \item \lstinline!desired_n_neighbors!: integer desired neighborhood size (default \lstinline!0!).
        \end{itemize}

        \textbf{Returns:}
        \begin{itemize}
            \item \lstinline!list!: \{
                \lstinline!neighborhood_residuals!: numeric array,\\
                \lstinline!neighborhood_indices!: integer matrix of gene indices
            \}
        \end{itemize}

        \begin{lstlisting}[language=R, caption={Neighborhood Construction}]
out <- tox_construct_neighborhoods(
  x_star,
  n_pool,
  mean_S,
  resid_S,
  desired_n_neighbors = 0
)

neigh_res <- out$neighborhood_residuals
neigh_idx <- out$neighborhood_indices
        \end{lstlisting}

    \item \textbf{Determine Shared Residual Range (Expert)}

        Compute the shared residual range \(R\) from a precomputed residual pool and its sorting permutation.  
        This expert interface assumes that residuals have already been flattened, cleaned, and sorted externally.

        \textbf{Arguments:}
        \begin{itemize}
            \item \lstinline!residual_pool!: numeric vector of absolute residuals (NaNs removed).
            \item \lstinline!residual_pool_perm!: integer vector giving the permutation that sorts \lstinline!residual_pool!.
            \item \lstinline!residual_range_quantile!: numeric scalar specifying the percentile (default \lstinline!95.0!).
        \end{itemize}

        \textbf{Returns:}
        \begin{itemize}
            \item \lstinline!numeric!: shared residual range \(R\).
        \end{itemize}

        \begin{lstlisting}[language=R, caption={Shared Residual Range (Expert)}]
R <- tox_determine_shared_residual_range_expert(
  residual_pool,
  residual_pool_perm,
  residual_range_quantile = 95.0
)
        \end{lstlisting}

    \item \textbf{Determine Shared Residual Range}

        Compute the shared residual range \(R\) from two residual matrices.  
        Both matrices are flattened, pooled, and the percentile threshold is applied to determine the symmetric range \([-R, R]\).

        \textbf{Arguments:}
        \begin{itemize}
            \item \lstinline!neighborhood_residuals_S1!: numeric matrix \lstinline!(n_residuals * n_neighbors)!.
            \item \lstinline!neighborhood_residuals_S2!: numeric matrix \lstinline!(n_residuals * n_neighbors)!.
            \item \lstinline!residual_range_quantile!: numeric scalar (default \lstinline!95.0!).
        \end{itemize}

        \textbf{Returns:}
        \begin{itemize}
            \item \lstinline!numeric!: shared residual range \(R\).
        \end{itemize}

        \begin{lstlisting}[language=R, caption={Shared Residual Range}]
R <- tox_determine_shared_residual_range(
  neighborhood_residuals_S1,
  neighborhood_residuals_S2,
  residual_range_quantile = 95.0
)
        \end{lstlisting}

    \item \textbf{Build Residual Histograms}

Build histogram bin counts and normalized probability mass functions (PMFs) for each neighbor.  
Residuals are binned using a shared residual range and a fixed number of bins.

\textbf{Arguments:}
\begin{itemize}
    \item \lstinline!neighborhood_residuals!: numeric matrix \lstinline!(n_residuals * n_neighbors)!.
    \item \lstinline!shared_residual_range!: numeric scalar \(R\).
    \item \lstinline!n_bins!: integer number of histogram bins.
\end{itemize}

\textbf{Returns:}
\begin{itemize}
    \item \lstinline!list!: \{
        \lstinline!counts!: integer matrix \lstinline!(n_neighbors * n_bins)!,\\
        \lstinline!pmf!: numeric matrix \lstinline!(n_neighbors * n_bins)!,\\
        \lstinline!included_n_residuals!: integer vector \lstinline!(n_neighbors)!
    \}
\end{itemize}

\begin{lstlisting}[language=R, caption={Residual Histogram Construction}]
out <- tox_build_residual_histograms(
  neighborhood_residuals,
  shared_residual_range,
  n_bins
)

counts <- out$counts
pmf <- out$pmf
included <- out$included_n_residuals
\end{lstlisting}

    \item \textbf{Build Residual Histograms (Filtered)}

        Build histogram bin counts and normalized probability mass functions (PMFs) for each neighbor, using a shared residual range and a fixed number of bins.  
        A logical or integer mask selects which neighbors contribute to each reference point’s histogram.

        \textbf{Arguments:}
        \begin{itemize}
            \item \lstinline!neighborhood_residuals!: numeric matrix \lstinline!(n_residuals * n_neighbors)!.
            \item \lstinline!shared_residual_range!: numeric scalar \(R\).
            \item \lstinline!n_bins!: integer number of histogram bins.
            \item \lstinline!neighbor_mask!: logical or integer matrix \lstinline!(n_neighbors * n_points)!, where non‑zero indicates inclusion.
        \end{itemize}

        \textbf{Returns:}
        \begin{itemize}
            \item \lstinline!list!: \{
                \lstinline!counts!: integer matrix \lstinline!(n_neighbors * n_bins)!,\\
                \lstinline!pmf!: numeric matrix \lstinline!(n_neighbors * n_bins)!,\\
                \lstinline!included_n_residuals!: integer vector \lstinline!(n_neighbors)!
            \}
        \end{itemize}

\begin{lstlisting}[language=R, caption={Filtered Residual Histogram Construction}]
out <- tox_build_residual_histograms_filtered(
  neighborhood_residuals,
  shared_residual_range,
  n_bins,
  neighbor_mask
)

counts   <- out$counts
pmf      <- out$pmf
included <- out$included_n_residuals
\end{lstlisting}

    \item \textbf{Compute Divergence per Reference Point}

        Compute the Jensen–Shannon divergence for each neighbor using two PMF matrices.

        \textbf{Arguments:}
        \begin{itemize}
            \item \lstinline!pmf_S1!: numeric matrix \lstinline!(n_neighbors * n_bins)!.
            \item \lstinline!pmf_S2!: numeric matrix \lstinline!(n_neighbors * n_bins)!.
        \end{itemize}

        \textbf{Returns:}
        \begin{itemize}
            \item \lstinline!numeric!: vector of length \lstinline!n_neighbors! containing JSD values.
        \end{itemize}

        \begin{lstlisting}[language=R, caption={Per‑Neighbor JSD}]
jsd <- tox_compute_divergence_per_reference_point(
  pmf_S1,
  pmf_S2
)
        \end{lstlisting}

    \item \textbf{Compute Weighted Global Divergence}

        Compute the weighted global Jensen–Shannon divergence from per‑neighbor divergences and included residual counts.

        \textbf{Arguments:}
        \begin{itemize}
            \item \lstinline!js_divergences!: numeric vector of length \lstinline!n_neighbors!.
            \item \lstinline!included_n_residuals_S1!: integer vector of length \lstinline!n_neighbors!.
            \item \lstinline!included_n_residuals_S2!: integer vector of length \lstinline!n_neighbors!.
        \end{itemize}

        \textbf{Returns:}
        \begin{itemize}
            \item \lstinline!list!: \{
                \lstinline!global_js_divergence!: numeric scalar,\\
                \lstinline!weights!: numeric vector \lstinline!(n_neighbors)!
            \}
        \end{itemize}

        \begin{lstlisting}[language=R, caption={Weighted Global JSD}]
out <- tox_compute_weighted_global_divergence(
  js_divergences,
  included_n_residuals_S1,
  included_n_residuals_S2
)

global_jsd <- out$global_js_divergence
weights <- out$weights
        \end{lstlisting}

    \item \textbf{Permutation Test for Global JSD}

        Estimate how likely the observed global Jensen–Shannon divergence is to occur by chance under the null hypothesis that both studies are exchangeable.  
        Residuals for each reference point are concatenated, shuffled, and reassigned to generate null‑distribution JSD values.

        \textbf{Arguments:}
        \begin{itemize}
            \item \lstinline!neighborhood_residuals_S1!: numeric vector or array of residuals for study 1.
            \item \lstinline!neighborhood_residuals_S2!: numeric vector or array of residuals for study 2.
            \item \lstinline!global_jsd_observed!: numeric scalar, the observed global weighted JSD.
            \item \lstinline!n_bins!: integer number of histogram bins.
            \item \lstinline!shared_residual_range!: numeric scalar specifying the shared residual range.
            \item \lstinline!n_permutations!: integer number of permutations.
            \item \lstinline!random_seed!: integer seed for reproducibility.
        \end{itemize}

        \textbf{Returns:}
        \begin{itemize}
            \item \lstinline!list!: \{
                \lstinline!jsd_null!: numeric vector of length \lstinline!n_permutations!,\\
                \lstinline!p_value!: numeric scalar
            \}
        \end{itemize}

        \begin{lstlisting}[language=R, caption={Permutation Test for Global JSD}]
out <- tox_gjct_permutation_test(
  neighborhood_residuals_S1,
  neighborhood_residuals_S2,
  global_jsd_observed,
  n_bins,
  shared_residual_range,
  n_permutations,
  random_seed = 42
)

jsd_null <- out$jsd_null
p_value  <- out$p_value
        \end{lstlisting}

    \item \textbf{Permutation Test for Global JSD (Filtered)}

        Estimate how likely the observed global Jensen–Shannon divergence is to occur by chance under the null hypothesis that both studies are exchangeable.  
        Neighbor masks allow excluding specific neighbors for each reference point in each study during histogram construction.

        \textbf{Arguments:}
        \begin{itemize}
            \item \lstinline!neighborhood_residuals_S1!: numeric array \lstinline!(n_reps * n_neighbors * n_points)! for study 1.
            \item \lstinline!neighborhood_residuals_S2!: numeric array \lstinline!(n_reps * n_neighbors * n_points)! for study 2.
            \item \lstinline!global_jsd_observed!: numeric scalar, the observed global weighted JSD.
            \item \lstinline!n_bins!: integer number of histogram bins.
            \item \lstinline!shared_residual_range!: numeric scalar specifying the shared residual range.
            \item \lstinline!n_permutations!: integer number of permutations.
            \item \lstinline!neighbor_mask_S1!: logical or integer matrix \lstinline!(n_neighbors * n_points)! for study 1.
            \item \lstinline!neighbor_mask_S2!: logical or integer matrix \lstinline!(n_neighbors * n_points)! for study 2.
            \item \lstinline!random_seed!: integer seed for reproducibility.
        \end{itemize}

        \textbf{Returns:}
        \begin{itemize}
            \item \lstinline!list!: \{
                \lstinline!jsd_null!: numeric vector of length \lstinline!n_permutations!,\\
                \lstinline!p_value!: numeric scalar
            \}
        \end{itemize}

\begin{lstlisting}[language=R, caption={Filtered Permutation Test for Global JSD}]
out <- tox_gjct_permutation_test_filtered(
  neighborhood_residuals_S1,
  neighborhood_residuals_S2,
  global_jsd_observed,
  n_bins,
  shared_residual_range,
  n_permutations,
  neighbor_mask_S1,
  neighbor_mask_S2,
  random_seed = 42
)

jsd_null <- out$jsd_null
p_value  <- out$p_value
\end{lstlisting}

    \item \textbf{Compute Family-Level JSD}

        Compute the family-level Jensen–Shannon divergence for a given paralog family.  
        Residuals are filtered by family membership, histograms are constructed, per‑point JSD values are computed, and a weighted global divergence is returned.

        \textbf{Arguments:}
        \begin{itemize}
            \item \lstinline!family_idx!: integer index of the gene family.
            \item \lstinline!gene_to_family_S1!: integer vector of length \lstinline!n_genes_S1! mapping genes in study 1 to families.
            \item \lstinline!gene_to_family_S2!: integer vector of length \lstinline!n_genes_S2! mapping genes in study 2 to families.
            \item \lstinline!neighborhood_residuals_S1!: numeric array \lstinline!(n_reps_S1 * n_neighbors * n_points)!.
            \item \lstinline!neighborhood_residuals_S2!: numeric array \lstinline!(n_reps_S2 * n_neighbors * n_points)!.
            \item \lstinline!neighborhood_genes_S1!: integer matrix \lstinline!(n_neighbors * n_points)!.
            \item \lstinline!neighborhood_genes_S2!: integer matrix \lstinline!(n_neighbors * n_points)!.
            \item \lstinline!n_bins!: integer number of histogram bins.
            \item \lstinline!shared_residual_range!: numeric scalar specifying the shared residual range.
        \end{itemize}

        \textbf{Returns:}
        \begin{itemize}
            \item \lstinline!list!: \{
                \lstinline!js_divergences!: numeric vector \lstinline!(n_points)!,\\
                \lstinline!included_n_reps_S1!: integer vector \lstinline!(n_points)!,\\
                \lstinline!included_n_reps_S2!: integer vector \lstinline!(n_points)!,\\
                \lstinline!total_included_n_reps!: integer scalar,\\
                \lstinline!global_js_divergence!: numeric scalar,\\
                \lstinline!weights!: numeric vector \lstinline!(n_points)!,\\
                \lstinline!ierr!: integer error code
            \}
        \end{itemize}

        \begin{lstlisting}[language=R, caption={Family-Level JSD}]
out <- tox_fjct_compute_jsd(
  family_idx,
  gene_to_family_S1,
  gene_to_family_S2,
  neighborhood_residuals_S1,
  neighborhood_residuals_S2,
  neighborhood_genes_S1,
  neighborhood_genes_S2,
  n_bins,
  shared_residual_range
)

jsd        <- out$js_divergences
global_jsd <- out$global_js_divergence
weights    <- out$weights
        \end{lstlisting}

    \item \textbf{Compute Family-Level JSD (Expert)}

        Compute family-level Jensen–Shannon divergence using precomputed neighbor masks.  
        This expert interface bypasses family‑mapping logic and directly applies masks to select neighbors for each reference point.

        \textbf{Arguments:}
        \begin{itemize}
            \item \lstinline!neighborhood_residuals_S1!: numeric array \lstinline!(n_reps_S1 * n_neighbors * n_points)!.
            \item \lstinline!neighborhood_residuals_S2!: numeric array \lstinline!(n_reps_S2 * n_neighbors * n_points)!.
            \item \lstinline!neighbor_mask_S1!: logical or integer matrix \lstinline!(n_neighbors * n_points)!.
            \item \lstinline!neighbor_mask_S2!: logical or integer matrix \lstinline!(n_neighbors * n_points)!.
            \item \lstinline!n_bins!: integer number of histogram bins.
            \item \lstinline!shared_residual_range!: numeric scalar specifying the shared residual range.
        \end{itemize}

        \textbf{Returns:}
        \begin{itemize}
            \item \lstinline!list!: \{
                \lstinline!js_divergences!,\\
                \lstinline!included_n_reps_S1!,\\
                \lstinline!included_n_reps_S2!,\\
                \lstinline!total_included_n_reps!,\\
                \lstinline!global_js_divergence!,\\
                \lstinline!weights!,\\
                \lstinline!pmf_S1!,\\
                \lstinline!pmf_S2!,\\
                \lstinline!tmp_counts!,\\
                \lstinline!ierr!
            \}
        \end{itemize}

        \begin{lstlisting}[language=R, caption={Family-Level JSD (Expert Interface)}]
out <- tox_fjct_compute_jsd_expert(
  neighborhood_residuals_S1,
  neighborhood_residuals_S2,
  neighbor_mask_S1,
  neighbor_mask_S2,
  n_bins,
  shared_residual_range
)

global_jsd <- out$global_js_divergence
weights    <- out$weights
        \end{lstlisting}

    \item \textbf{Compute Contribution Scores}

        Compute per‑family contribution scores from global JSD values and the number of included residuals per family.  
        Contribution scores quantify how strongly each family contributes to the overall divergence.

        \textbf{Arguments:}
        \begin{itemize}
            \item \lstinline!global_js_divergences!: numeric vector of length \lstinline!k_families!.
            \item \lstinline!total_included_n_reps_per_f!: integer vector of length \lstinline!k_families!.
        \end{itemize}

        \textbf{Returns:}
        \begin{itemize}
            \item \lstinline!list!: \{
                \lstinline!support_weights!: numeric vector,\\
                \lstinline!contribution_scores!: numeric vector,\\
                \lstinline!ierr!: integer error code
            \}
        \end{itemize}

        \begin{lstlisting}[language=R, caption={Contribution Score Computation}]
out <- tox_fjct_compute_contribution_scores(
  global_js_divergences,
  total_included_n_reps_per_f
)

scores <- out$contribution_scores
weights <- out$support_weights
        \end{lstlisting}
\end{enumerate}
